\documentclass[12pt, a4paper]{ctexart}
\usepackage{geometry}
\geometry{left=2.5cm, right=2.5cm, top=3cm, bottom=3cm}
\usepackage{listings}
\usepackage{graphicx}
\usepackage{xcolor}
\usepackage{amsmath}
\usepackage{booktabs}
\usepackage{tikz}
\usetikzlibrary{positioning, arrows.meta}

% Code snippet settings
\lstset{
    basicstyle=\ttfamily\small,
    frame=single,
    framesep=5pt,
    xleftmargin=10pt,
    xrightmargin=10pt,
    keywordstyle=\color{blue},
    commentstyle=\color{gray},
    breaklines=true
}

\begin{document}

% 封面
\begin{titlepage}
    \centering
    \vspace*{5cm}
    {\Huge\bfseries 实验报告\par}
    \vspace{2cm}
    {\Large 课程：数字逻辑与计算机组成实验\par}
    \vspace{1cm}
    {\Large 实验十一：字符交互界面\par}
    \vspace{4cm}
    {\large 姓名：\underline{\makebox[4cm]{郑袭明}}\par}
    \vspace{1cm}
    {\large 学号：\underline{\makebox[4cm]{251502027}}\par}
    \vspace{1cm}
    {\large 日期：\today\par}
\end{titlepage}

\tableofcontents
\newpage

\section{设计思路与架构差异}

本实验实现了一个可以通过键盘输入并在 VGA 显示器上回显的字符交互界面。相比于实验指导书的推荐设计，本实现进行了多项优化，
并完成了所有可选的扩展要求（如自动滚屏、Backspace 跨行删除、Shift 键大小写与特殊符号切换、长按键重复输入等）。

\subsection{字符显示分辨率与地址计算简化}

实验指导书要求采用宽 9 像素、高 16 像素的字模，并在 $640 \times 480$ 的屏幕上显示 $70 \times 30$ 的字符阵列。
但在 FPGA 中，对 9 像素宽度的字符进行寻址涉及除以 9 或模 9 的操作，这在硬件上会消耗大量乘除法资源，并引入极长的组合逻辑延迟，严重影响时序收敛。

为此，本实现做出了如下优化：
\begin{enumerate}
    \item \textbf{字符尺寸调整}：将字符宽度改为 8 像素，高度保持为 16 像素。此时屏幕分辨率实际为 $80 \times 30$ 个字符。
    \item \textbf{字模兼容处理}：直接使用 12 位宽的 \texttt{font\_rom}，但在显示提取像素时，仅通过 \texttt{h\_addr[2:0]}索引低 8 位，自动舍弃了第 9 位（为字符间距的空列）。
    \item \textbf{无乘除法寻址}：利用 2 的幂次宽度，列地址计算简化为简单的位拼接与移位操作。对于 VGA 控制器输出的当前像素坐标 $(h\_addr, v\_addr)$，其对应的字符 VRAM 单元地址计算为：
    \begin{equation}
        vram\_addr = (vram\_row \times 80) + h\_addr[9:3]
    \end{equation}
    其中 $\times 80$ 可由 $(vram\_row \ll 6) + (vram\_row \ll 4)$ 的无符号移位加法电路高效实现，规避了乘法器，极大优化了时序性能。
\end{enumerate}

\subsection{硬件高效环形滚屏设计}

为了支持“输入至最后一行后自动换行，且已输入行自动上移”的滚屏需求，传统的软件做法是整体平移 VRAM 中的数据。但这在硬件上需要一次性拷贝 $30 \times 80 = 2400$ 字节，极其消耗时钟周期与逻辑资源。

本实现采用\textbf{硬件级环形缓冲区（Circular Buffer）}实现了 $O(1)$ 的零拷贝自动滚屏：
\begin{itemize}
    \item 引入 1 位寄存器 \texttt{offset} 记录是否发生过滚动，以及 \texttt{cursor\_y} 记录当前光标所在的虚拟 VRAM 行。
    \item 当光标从最后一行（第 29 行）换行时，不平移显存，而是将 \texttt{offset} 置 1，并将光标坐标 \texttt{cursor\_y} 环形累加。同时，在写入下一行前，将下一行的 VRAM 有效位一次性清空：
    \begin{equation}
        vram\_valid[next\_y \times 80 +: 80] \leftarrow 0
    \end{equation}
    \item 在 VGA 扫描输出时，根据当前扫描的物理行 \texttt{v\_addr[9:4]} 与光标位置 \texttt{cursor\_y} 进行动态地址重映射：
    \begin{align}
        y_{sum} &= v\_addr[9:4] + cursor\_y + 1 \\
        vram\_row &= \begin{cases}
            y_{sum} - 30 & y_{sum} \ge 30 \\
            y_{sum}[4:0] & y_{sum} < 30
        \end{cases}
    \end{align}
    这样，屏幕顶行（物理第 0 行）总是映射显示虚拟的 \texttt{(cursor\_y + 1) \% 30} 行，屏幕底行总是映射显示当前光标所在的 \texttt{cursor\_y} 行，从而在零内存拷贝的开销下实现了平滑的硬件滚屏。
\end{itemize}

\subsection{显存读写时序与冲突处理}

字符显存（VRAM）同时被 VGA 扫描逻辑高频读取，又被键盘输入逻辑随机写入。
\begin{itemize}
    \item \textbf{异步/组合读}：由于 VGA 像素时钟为 25MHz，读取显存及字模 ROM 采用了纯组合逻辑路径：
    \begin{center}
        \texttt{vram\_addr} $\to$ \texttt{char\_code} $\to$ \texttt{font\_row} $\to$ \texttt{vga\_data}
    \end{center}
    此方式能够零延迟响应 VGA 扫描信号，避免了引入同步读带来的时钟延迟偏斜和像素对齐问题。
    \item \textbf{同步写}：键盘模块的数据写入同步于 100MHz 的系统主时钟 \texttt{clk}。由于键盘按键的写入频率极低（人类极限输入速度折合几十赫兹），
    虽然存在读写时钟域的跨越，但实际碰撞概率极小，且即使发生碰撞也仅在极短的一帧内出现某像素的轻微闪烁，不会对视觉造成任何不良影响。
    这种“异步组合读 + 同步写”的设计在保证显示时序收敛的前提下，极大简化了冲突处理逻辑。
\end{itemize}

\subsection{键盘输入与特殊功能支持}

本实现在 \texttt{keyboard.v} 和 \texttt{cli.v} 中拓展并实现了以下特性：
\begin{enumerate}
    \item \textbf{Shift 键与符号映射}：通过在键盘状态机中监测左/右 Shift 键的 \texttt{make/break} 码，维护 \texttt{shift} 状态信号。
    在输出 ASCII 时，利用组合逻辑对小写字母进行大小写转换（$- 8'h20$），并对数字键和符号键进行重映射（例如 \texttt{8'h31} ($1$) 映射为 \texttt{8'h21} ($!$)），完整支持了全键盘符号输入。
    \item \textbf{跨行 Backspace 删除}：在 \texttt{cli.v} 中，当按下退格键（\texttt{8'h08}）且光标处于行首（\texttt{cursor\_x == 0}）时，若当前不是屏幕顶行，
    光标将自动回退到上一行的末尾（\texttt{cursor\_x <= 79}，\texttt{cursor\_y <= prev\_y}），并清除该位置的字符有效位。这实现了跨行删除的连贯交互。
    \item \textbf{硬件按键重复（Key Repeat）}：通过捕获 PS/2 键盘连续发送的相同按键 \texttt{make} 码（不伴随 \texttt{break} 码），
    在每次接收到新扫描码时累加事件计数器 \texttt{key\_count}。主控模块 \texttt{cli} 仅需检测 \texttt{key\_count} 的变化即可执行输入，
    直接利用了 PS/2 硬件协议的自动重发机制，免去了在 FPGA 内部设计复杂定时器的繁琐。
\end{enumerate}

\section{实验总结}

实验时只遇到了一个问题，\texttt{scancode\_ram.b} 是之前实验提供的，神奇的使用了同步的写，导致了一系列问题；改成异步就都修好了。

主要的时间就花在调试滚屏上了。

\end{document}
